\section{Further incidents and excerpt provenance}
\label{app:incidents}

The seven incidents in the main text were chosen for the mechanism each shows. The cases below are retained because they limit how the main incidents should be read.

\AppIncidentTitle{g11}{G11: a false completion claim, a correction, and two refusals}

Late in game 11 only Grok and Sol remain, and Grok holds a developed dark-blue group. Sol proposes:

\TranscriptQuote{fx-g11-offer}

Grok holds one pink deed and Sol holds the other two, so Virginia alone does not complete the group. Grok's private note calls the offer a lie. Sol's own note on the proposal repeats the false premise, which makes a mistake more likely than a bluff. Grok answers:

\TranscriptQuote{fx-g11-reject}

Sol concedes and offers a true completion:

\TranscriptQuote{fx-g11-concede}

\TranscriptQuote{fx-g11-reject-again}

Grok's note on the second refusal says the pink group is weak next to the position it already holds and that Pennsylvania would complete Sol's greens. This is a second refusal of one's own completion on rival-sensitive grounds, in an easier setting than Incident~\ref{inc:swap} because the refuser is far ahead. It also shows a recipient's accusation of lying that the speaker's own note does not support. All four are standard first attempts at G11 d255--258, turn 97.

\AppIncidentTitle{g3}{G3: a stated obstacle contradicted by the speaker's note}

Late in game 3 one house remains in the bank and Gemini could add a house to Park Place. Grok says in public that this needs someone to free buildings:

\TranscriptQuote{G3-2}

\TranscriptQuote{G3-3}

Gemini's next note gives its own reason for holding:

\TranscriptQuote{G3-5}

In the retained exact-prefix check, Gemini could have mortgaged two deeds and bought the house with the one remaining in the bank. The public claim is false for that route, and the speaker's note records the route. Gemini had already held at its previous decision, before the claim, and its note cites income. Exposure followed by a hold does not show reliance.

\AppIncidentTitle{g20leak}{G20: a private note delivered as public chat}

On Grok's first end-of-turn decision in game 20 the public-message argument of the received tool call contains a fragment of the private note:

\TranscriptQuote{g20-received-and-applied-public}

The engine publishes it, and three decisions later an opponent's note records it:

\TranscriptQuote{g20-recipient-recognition}

The assistant text of the same response requested an empty public message. The leak is real and was noticed. Its content concerns public holdings, and nothing shows it was exploited or intended.

\paragraph{Shorter notes.} In game 4, Grok states a false \$200 house cost for the yellow group and Gemini replies with a public concession, but Gemini's note one decision earlier records the true \$150 and the fallback it then takes. In game 6, Sol's note rejects a false claim about an orange monopoly and Sol accepts the trade anyway. In game 18, a Grok tool call accepts \$130 for Baltic Avenue while its text rejects, which completes the eventual winner's first color group. In game 12, a speaker calls its earlier insolvency pitch a bluff only after the recipient corrects the arithmetic, and its original explanation treats the figure as real, so the pitch is not counted as a knowing misstatement. These four are inherited from the earlier source reviews.

\subsection{Selection and provenance of excerpts}
\label{app:quotes}

Every card is a verbatim, contiguous span of one named field. Square-bracketed ellipses join two spans of the same field. Spelling and arithmetic inside excerpts are untouched. The card header gives the speaker, the channel, the game and zero-based decision, and the zero-based index of the selected attempt. A public message is the text emitted with the applied action. An applied private note is the \texttt{private\_thought} field of the applied action. A returned provider explanation is reasoning text in the provider response. Assistant content is the text part of a response whose tool call was applied.

Twenty-eight cards come from the retained transcript manifest of 29 October. For this draft a script reloaded each of their 28 source fields from the canonical saved-game files, compared file and field hashes with the manifest, and checked all 30 spans and their typeset form. Fourteen quotes are new in this draft: eleven cards for games 15, 1 and 11 and three inline quotes for game 16. The script extracts them from the resolved decision records and fails unless each requested span occurs exactly once in the field. Their offsets, hashes, file paths and line numbers are in \artifact{data/transcript_manifest_supplement.json}. The surrounding facts for the new incidents come from the offer and flow ledgers. They have not had the separate bounded source review that the 29 October incidents received, and recipient-prompt delivery was not re-checked for them beyond the engine's delivery of the offer itself.

Six of the seven main incidents involve Grok as a speaker. Its notes are long and specific, and it is the only system whose assistant text can be compared with its tool call. Opus's notes are mostly absent. The selection therefore reflects what can be observed, and it does not rank the systems by conduct.
